For $n=4$, the values are $24,26,27$, with $26=2\cdot13$ and $27=3\cdot9$. In general, for $2\le k\le n$, the integer $k$ divides $n!$ and therefore divides $n!+k$. Moreover $1<k<n!+k$, so this is a proper divisor and the number is composite. Given any positive integer $L$, take $n=L+1$. The $L$ consecutive integers $n!+2$ through $n!+n$ are all composite. Thus composite runs have unbounded length. This argument does not assert anything about infinitely many prime pairs with gap 2.
